% ---- score2 standard preamble (from: python3 score2_tools.py preamble) ----
\documentclass[10pt]{extarticle}
\usepackage[a4paper,margin=1.5cm]{geometry}
\usepackage[T1]{fontenc}
\usepackage{amsmath,amssymb}
\usepackage{graphicx}
\usepackage[export]{adjustbox}
\usepackage{multicol}
\usepackage{enumitem}
\usepackage{booktabs,array,tabularx}
\setlength{\columnsep}{0.7cm}
\setlength{\parindent}{0pt}
\setlength{\parskip}{3pt}
\pagestyle{empty}
\makeatletter
% figure registry (figdefs.tex is generated by the extract step)
\newcommand{\qfigdef}[4]{\@namedef{qf@file@#1}{#2}\@namedef{qf@w@#1}{#3}\@namedef{qf@wide@#1}{#4}}
\newcommand{\qf@check}[1]{\@ifundefined{qf@file@#1}{\PackageError{score2}{Unknown figure id #1}{}}{}}
\newcommand{\qf@img}[4][]{\includegraphics[width=#2\dimexpr\@nameuse{qf@w@#4}pt\relax,max width=#3,max height=0.45\textheight,#1]{\@nameuse{qf@file@#4}}}
% Full-width questions: finalize wraps every question that contains a wide figure in \qfullbegin...\qfullend,
% so the WHOLE question (stem, figures, options) is set across the page in reading order, then the two columns
% resume. The question's head - its number, stem and first figure (\qfullsplit marks the end of the head) -
% is kept on one page, so a number is never left at a page foot with its figure overleaf; the rest may break
% across pages normally, so no big gaps are left. Without a \qfullsplit, the whole question is kept together
% if it is under 45% of a page. Never write these yourself.
\newif\ifqf@full
\newif\ifqf@boxed
\newsavebox{\qf@box}
\newcommand{\qf@open}{\global\qf@fulltrue\par\global\qf@boxedtrue
  \setbox\qf@box\vbox\bgroup\hsize\textwidth\linewidth\textwidth}
\newcommand{\qf@emit}[1]{\par\egroup\global\qf@boxedfalse
  \ifdim\dimexpr\ht\qf@box+\dp\qf@box\relax<#1\textheight\noindent\usebox{\qf@box}\par\else\unvbox\qf@box\fi}
\newcommand{\qfullsplit}{\ifqf@boxed\qf@emit{0.85}\fi}
\newcommand{\qf@close}{\ifqf@boxed\qf@emit{0.45}\fi\global\qf@fullfalse}
\newcommand{\qfullbegin}{\end{multicols}\qf@open}
\newcommand{\qfullend}{\qf@close\begin{multicols}{2}}
\newcommand{\qfullstart}{\qf@open}
\newcommand{\qfulldone}{\qf@close}
% \qfig[scale]{id}: block figure, centred, directly where it is written (size set by the tool).
% Fallback if a wide figure is not inside a full-width question: break out of the columns in place - no floats.
\newcommand{\qfig}[2][1]{\qf@check{#2}%
  \ifqf@full\qf@blk{#1}{#2}\else
  \ifnum\@nameuse{qf@wide@#2}=1 \expandafter\@firstoftwo\else\expandafter\@secondoftwo\fi
  {\end{multicols}\begin{center}\qf@img{#1}{\textwidth}{#2}\end{center}\begin{multicols}{2}}%
  {\qf@blk{#1}{#2}}\fi}
\newcommand{\qf@blk}[2]{\par\smallskip{\centering\qf@img{#1}{\linewidth}{#2}\par}\smallskip}
% \qfigi[scale]{id}: inline figure (inside option lists / tables / minipages); never breaks columns.
\newcommand{\qfigi}[2][1]{\qf@check{#2}\qf@img[valign=c]{#1}{\linewidth}{#2}}
\makeatother
% structure macros - lint reads these to check completeness and figure placement
\newcommand{\papersec}[1]{\par\bigskip{\centering\Large\bfseries PAPER-#1\par}\medskip}
\newcommand{\partsec}[2]{\par\medskip\noindent{\large\bfseries PART-#1: #2}\par\smallskip}
\newcommand{\sect}[1]{\par\smallskip\noindent\textbf{#1}\par}
\newcommand{\qn}[1]{\par\smallskip\noindent\textbf{#1.}~}
\newcommand{\paraq}[2]{\par\smallskip\noindent\textbf{Paragraph for Questions #1 and #2}\par}
\newenvironment{qoptions}{\begin{itemize}[leftmargin=1.8em,itemsep=1pt,topsep=2pt,label={}]}{\end{itemize}}
\newcommand{\opt}[2]{\item[(#1)] #2}
% figure registry (generated by score2_tools.py finalize)
\qfigdef{p1_phy_s1_q01_1}{TOAS - Test 07 (01 May 2022)-img/p1_phy_s1_q01_1.png}{116.3}{0}
\qfigdef{p1_phy_s1_q02_B}{TOAS - Test 07 (01 May 2022)-img/p1_phy_s1_q02_B.png}{92.0}{0}
\qfigdef{p1_phy_s1_q03_1}{TOAS - Test 07 (01 May 2022)-img/p1_phy_s1_q03_1.png}{114.4}{0}
\qfigdef{p1_phy_s2_q05_1}{TOAS - Test 07 (01 May 2022)-img/p1_phy_s2_q05_1.png}{119.9}{0}
\qfigdef{p1_phy_s2_q06_1}{TOAS - Test 07 (01 May 2022)-img/p1_phy_s2_q06_1.png}{106.0}{0}
\qfigdef{p1_phy_s2_q07_1}{TOAS - Test 07 (01 May 2022)-img/p1_phy_s2_q07_1.png}{104.8}{0}
\qfigdef{p1_phy_s2_q10_1}{TOAS - Test 07 (01 May 2022)-img/p1_phy_s2_q10_1.png}{119.5}{0}
\qfigdef{p1_phy_s3_para1_2_1}{TOAS - Test 07 (01 May 2022)-img/p1_phy_s3_para1_2_1.png}{106.1}{0}
\qfigdef{p1_phy_s3_para3_4_1}{TOAS - Test 07 (01 May 2022)-img/p1_phy_s3_para3_4_1.png}{72.3}{0}
\qfigdef{p1_phy_s3_para5_6_1}{TOAS - Test 07 (01 May 2022)-img/p1_phy_s3_para5_6_1.png}{103.8}{0}
\qfigdef{p1_phy_s4_q08_1}{TOAS - Test 07 (01 May 2022)-img/p1_phy_s4_q08_1.png}{107.8}{0}
\qfigdef{p1_chem_s1_q01_1}{TOAS - Test 07 (01 May 2022)-img/p1_chem_s1_q01_1.png}{141.1}{0}
\qfigdef{p1_chem_s1_q01_A}{TOAS - Test 07 (01 May 2022)-img/p1_chem_s1_q01_A.png}{33.0}{0}
\qfigdef{p1_chem_s1_q01_B}{TOAS - Test 07 (01 May 2022)-img/p1_chem_s1_q01_B.png}{30.1}{0}
\qfigdef{p1_chem_s1_q01_C}{TOAS - Test 07 (01 May 2022)-img/p1_chem_s1_q01_C.png}{34.2}{0}
\qfigdef{p1_chem_s1_q01_D}{TOAS - Test 07 (01 May 2022)-img/p1_chem_s1_q01_D.png}{33.9}{0}
\qfigdef{p1_chem_s1_q02_A}{TOAS - Test 07 (01 May 2022)-img/p1_chem_s1_q02_A.png}{132.7}{0}
\qfigdef{p1_chem_s1_q02_B}{TOAS - Test 07 (01 May 2022)-img/p1_chem_s1_q02_B.png}{158.9}{0}
\qfigdef{p1_chem_s1_q02_C}{TOAS - Test 07 (01 May 2022)-img/p1_chem_s1_q02_C.png}{128.1}{0}
\qfigdef{p1_chem_s2_q05_1}{TOAS - Test 07 (01 May 2022)-img/p1_chem_s2_q05_1.png}{99.9}{0}
\qfigdef{p1_chem_s2_q05_A}{TOAS - Test 07 (01 May 2022)-img/p1_chem_s2_q05_A.png}{38.7}{0}
\qfigdef{p1_chem_s2_q05_B}{TOAS - Test 07 (01 May 2022)-img/p1_chem_s2_q05_B.png}{46.0}{0}
\qfigdef{p1_chem_s2_q05_C}{TOAS - Test 07 (01 May 2022)-img/p1_chem_s2_q05_C.png}{39.2}{0}
\qfigdef{p1_chem_s2_q05_D}{TOAS - Test 07 (01 May 2022)-img/p1_chem_s2_q05_D.png}{46.1}{0}
\qfigdef{p1_chem_s3_para1_2_1}{TOAS - Test 07 (01 May 2022)-img/p1_chem_s3_para1_2_1.png}{86.5}{0}
\qfigdef{p1_chem_s3_q04_1}{TOAS - Test 07 (01 May 2022)-img/p1_chem_s3_q04_1.png}{245.1}{0}
\qfigdef{p1_chem_s4_q07_1}{TOAS - Test 07 (01 May 2022)-img/p1_chem_s4_q07_1.png}{332.1}{1}
\qfigdef{p2_phy_s1_q01_1}{TOAS - Test 07 (01 May 2022)-img/p2_phy_s1_q01_1.png}{44.3}{0}
\qfigdef{p2_phy_s1_q02_1}{TOAS - Test 07 (01 May 2022)-img/p2_phy_s1_q02_1.png}{115.9}{0}
\qfigdef{p2_phy_s1_q04_1}{TOAS - Test 07 (01 May 2022)-img/p2_phy_s1_q04_1.png}{54.0}{0}
\qfigdef{p2_phy_s1_q05_1}{TOAS - Test 07 (01 May 2022)-img/p2_phy_s1_q05_1.png}{94.6}{0}
\qfigdef{p2_phy_s2_para7_8_1}{TOAS - Test 07 (01 May 2022)-img/p2_phy_s2_para7_8_1.png}{128.3}{0}
\qfigdef{p2_phy_s2_para9_10_1}{TOAS - Test 07 (01 May 2022)-img/p2_phy_s2_para9_10_1.png}{144.0}{0}
\qfigdef{p2_phy_s2_q09_A}{TOAS - Test 07 (01 May 2022)-img/p2_phy_s2_q09_A.png}{107.0}{0}
\qfigdef{p2_phy_s2_q09_B}{TOAS - Test 07 (01 May 2022)-img/p2_phy_s2_q09_B.png}{104.4}{0}
\qfigdef{p2_phy_s2_q09_C}{TOAS - Test 07 (01 May 2022)-img/p2_phy_s2_q09_C.png}{105.5}{0}
\qfigdef{p2_phy_s2_q09_D}{TOAS - Test 07 (01 May 2022)-img/p2_phy_s2_q09_D.png}{105.8}{0}
\qfigdef{p2_phy_s3_q02_1}{TOAS - Test 07 (01 May 2022)-img/p2_phy_s3_q02_1.png}{57.4}{0}
\qfigdef{p2_phy_s3_para3_4_1}{TOAS - Test 07 (01 May 2022)-img/p2_phy_s3_para3_4_1.png}{109.4}{0}
\qfigdef{p2_phy_s3_para5_6_1}{TOAS - Test 07 (01 May 2022)-img/p2_phy_s3_para5_6_1.png}{235.6}{0}
\qfigdef{p2_phy_s4_q07_1}{TOAS - Test 07 (01 May 2022)-img/p2_phy_s4_q07_1.png}{154.3}{0}
\qfigdef{p2_phy_s4_q08_1}{TOAS - Test 07 (01 May 2022)-img/p2_phy_s4_q08_1.png}{67.3}{0}
\qfigdef{p2_phy_s4_q09_1}{TOAS - Test 07 (01 May 2022)-img/p2_phy_s4_q09_1.png}{119.6}{0}
\qfigdef{p2_chem_s1_q01_A}{TOAS - Test 07 (01 May 2022)-img/p2_chem_s1_q01_A.png}{127.1}{0}
\qfigdef{p2_chem_s1_q01_2}{TOAS - Test 07 (01 May 2022)-img/p2_chem_s1_q01_2.png}{124.9}{0}
\qfigdef{p2_chem_s1_q01_C}{TOAS - Test 07 (01 May 2022)-img/p2_chem_s1_q01_C.png}{123.9}{0}
\qfigdef{p2_chem_s1_q01_4}{TOAS - Test 07 (01 May 2022)-img/p2_chem_s1_q01_4.png}{50.2}{0}
\qfigdef{p2_chem_s1_q02_1}{TOAS - Test 07 (01 May 2022)-img/p2_chem_s1_q02_1.png}{154.2}{0}
\qfigdef{p2_chem_s2_para7_8_1}{TOAS - Test 07 (01 May 2022)-img/p2_chem_s2_para7_8_1.png}{292.9}{1}
\qfigdef{p2_chem_s2_q07_1}{TOAS - Test 07 (01 May 2022)-img/p2_chem_s2_q07_1.png}{113.4}{0}
\qfigdef{p2_chem_s2_q07_A}{TOAS - Test 07 (01 May 2022)-img/p2_chem_s2_q07_A.png}{48.5}{0}
\qfigdef{p2_chem_s2_q07_B}{TOAS - Test 07 (01 May 2022)-img/p2_chem_s2_q07_B.png}{50.6}{0}
\qfigdef{p2_chem_s2_q07_C}{TOAS - Test 07 (01 May 2022)-img/p2_chem_s2_q07_C.png}{48.7}{0}
\qfigdef{p2_chem_s2_q07_D}{TOAS - Test 07 (01 May 2022)-img/p2_chem_s2_q07_D.png}{51.9}{0}
\qfigdef{p2_chem_s2_q08_1}{TOAS - Test 07 (01 May 2022)-img/p2_chem_s2_q08_1.png}{143.4}{0}
\qfigdef{p2_chem_s3_para3_4_1}{TOAS - Test 07 (01 May 2022)-img/p2_chem_s3_para3_4_1.png}{123.5}{0}
\raggedcolumns
\makeatletter\@beginparpenalty=10000\makeatother
\newcolumntype{L}[1]{>{\raggedright\arraybackslash}p{#1}}
\begin{document}
\begin{center}{\Large\bfseries ALLEN TOAS --- Test 07 (01 May 2022)}\end{center}
\begin{multicols}{2}
\raggedcolumns
\papersec{1}
\partsec{1}{PHYSICS}
\sect{SECTION-I(i)}
\qn{1} A cart loaded with sand moves along a horizontal floor due to a constant force $F$ coinciding in direction with the cart's velocity vector. In the process sand spills through a hole. in the bottom with constant rate $\mu$ kg/sec. Find the velocity of the cart as a function of time `$t$', if at the initial moment $t=0$ the cart with loaded sand had the mass $m_0$ and its velocity was equal to zero.
\qfig{p1_phy_s1_q01_1}
\begin{qoptions}
\opt{A}{$\dfrac{F}{\mu}\ln\!\left(\dfrac{m_0}{m_0-\mu t}\right)$}
\opt{B}{$\dfrac{F}{\mu}\ln\!\left(\dfrac{m_0}{m_0+\mu t}\right)$}
\opt{C}{$\dfrac{F}{\mu}\ln\!\left(\dfrac{m_0+\mu t}{m_0}\right)$}
\opt{D}{$\dfrac{F}{\mu}\ln\!\left(\dfrac{m_0-\mu t}{m_0}\right)$}
\end{qoptions}
\qn{2} A fixed non-conducting spherical shell of radius $R$ has a charge $Q$ spread uniformly on its surface. A particle having mass `$m$' and charge `$q$' is projected towards the shell from a large distance at a speed $u$. Impact parameter is $\dfrac R2$ as shown in figure. The particle passes the shell grazing its surface. The projection speed $u$ of the particle is (Neglect gravity).
\qfig{p1_phy_s1_q02_B}
\begin{qoptions}
\opt{A}{$\sqrt{\dfrac{4Qq}{3\pi\epsilon_0 mR}}$}
\opt{B}{$\sqrt{\dfrac{Qq}{2\pi\epsilon_0 mR}}$}
\opt{C}{$\sqrt{\dfrac{2Qq}{3\pi\epsilon_0 mR}}$}
\opt{D}{$\sqrt{\dfrac{3Qq}{2\pi\epsilon_0 mR}}$}
\end{qoptions}
\qn{3} A small block of mass $m$ resting on a horizontal smooth floor at distance $\ell$ from a rigid wall. Block attached to the spring, is pushed towards right by a distance $\dfrac{3\ell}{2}$ and released. When block passes from its mean position an other block of mass $m_1$ is placed on it which sticks to it due to friction. If combined block just collides with the left wall, then :-
\qfig{p1_phy_s1_q03_1}
\begin{qoptions}
\opt{A}{$m_1=m$}
\opt{B}{$m_1=\dfrac m4$}
\opt{C}{$m_1=\dfrac{3m}{4}$}
\opt{D}{$m_1=\dfrac{5m}{4}$}
\end{qoptions}
\qn{4} A point object is placed at a distance of 15 cm from a convex lens. The image is formed on the other side of lens at a distance 30 cm from lens. When a concave lens is placed in contact with convex lens, image is shifted further by 30 cm. Calculate the focal length of concave lens :-
\begin{qoptions}
\opt{A}{10 cm}
\opt{B}{$-10$ cm}
\opt{C}{$-60$ cm}
\opt{D}{60 cm}
\end{qoptions}
\sect{SECTION-I(ii)}
\qn{5} Two blocks $A$ and $B$ each of mass `$m$', are connected by a massless spring of force constant $k$. Initially given system is at rest and spring is in its natural length. A constant force $F$ starts acting on block $A$ at time $t=0$ and at time $t$, the extension in the spring is seen to be $\ell$. Given system lies on smooth horizontal surface.
\qfig{p1_phy_s2_q05_1}
\begin{qoptions}
\opt{A}{Acceleration of centre of mass of the given system ($A+B+$ spring) is $\dfrac Fm$.}
\opt{B}{Displacement of centre of mass of the given system ($A+B+$ spring) in time $t$ is $\dfrac{Ft^2}{4m}$.}
\opt{C}{Displacement of block $A$ in time $t$ is $\dfrac{Ft^2}{4m}+\dfrac{\ell}{2}$.}
\opt{D}{Displacement of block $A$ in time $t$ is $\dfrac{Ft^2}{2m}-\dfrac{\ell}{2}$.}
\end{qoptions}
\qn{6} An infinite line charge is perpendicular to the plane of the figure having linear charge density $\lambda$. A particle having charge $+Q$ and mass $m$ is projected in the field of the line charge from point $P$. The point $P$ is at a distance $R$ from the line charge and velocity given to the particle is $v_0$ perpendicular to the radial line at $P$ as shown in figure.
\qfig{p1_phy_s2_q06_1}
\begin{qoptions}
\opt{A}{Speed of the particle when its distance from the line charge grows to $\eta R$ ($\eta>1$), is $\sqrt{v_0^2+\dfrac{\lambda Q}{\pi\epsilon_0 m}\ln(\eta)}$}
\opt{B}{Speed of the particle when its distance from the line charge grows to $\eta R$ ($\eta>1$), is $\sqrt{v_0^2-\dfrac{\lambda Q}{\pi\epsilon_0 m}\ln(\eta)}$}
\opt{C}{Velocity component of the particle along the radial line (joining the line charge to the particle) at the instant its distance is $\eta R$, is given by $\sqrt{v_0^2\left(1-\dfrac1{\eta^2}\right)+\dfrac{\lambda Q}{\pi\epsilon_0 m}\ln(\eta)}$}
\opt{D}{Velocity component of the particle along the radial line (joining the line charge to the particle) at the instant its distance is $\eta R$, is given by $\sqrt{v_0^2\left(1+\dfrac1{\eta^2}\right)+\dfrac{\lambda Q}{\pi\epsilon_0 m}\ln(\eta)}$}
\end{qoptions}
\qn{7} Figure shows an electric dipole formed by two particles fixed at the ends of a light rod of length $\ell$. The mass of each particle is $m$ and the charges are $-q$ and $+q$. The system is placed in such a way that the dipole axis is parallel to a uniform electric field $E$ that exist in the region. The dipole is slightly rotated about its centre and released. Considering small angular displacement $\theta$, which of the following options are correct.
\qfig{p1_phy_s2_q07_1}
\begin{qoptions}
\opt{A}{Motion of the dipole is angular SHM with time period $2\pi\sqrt{\dfrac{m\ell}{2qE}}$}
\opt{B}{Motion of the dipole is angular SHM with time period $2\pi\sqrt{\dfrac{m\ell}{qE}}$}
\opt{C}{During SHM no net force acts on dipole.}
\opt{D}{During SHM, when dipole moment vector becomes parallel to electric field, net torque on dipole is zero.}
\end{qoptions}
\qn{8} Two fixed charges $-2Q$ and $Q$ are located at the points with coordinates $(-3a,0)$ and $(+3a,0)$ respectively in the $xy$ plane.
\begin{qoptions}
\opt{A}{All the points in $x$-$y$ plane where the electric potential due to the two charges is zero lie on a circle of radius $4a$.}
\opt{B}{All the points in $x$-$y$ plane where the electric potential due to the two charges is zero lie on a circle having its centre at $(5a,0)$}
\opt{C}{All the points in $x$-$y$ plane where the electric potential due the two charges is zero lie on an ellipse.}
\opt{D}{All the points in $x$-$y$ plane where the electric potential due the two charges is zero lie on an ellipse having its centre at $(5a,0)$}
\end{qoptions}
\qn{9} A horizontal spring mass system of mass $M$ executes oscillatory motion of amplitude $a_0$ and time period $T_0$. When the mass $M$ is passing through its equilibrium position another mass $m$ is placed on it such that both move together. If $a$ and $T$ be the new amplitude and time period respectively then,
\begin{qoptions}
\opt{A}{$a=\sqrt{\dfrac{M}{M+m}}\,a_0$}
\opt{B}{$a=\sqrt{\dfrac{M+m}{M}}\,a_0$}
\opt{C}{$T=\sqrt{\dfrac{M}{M+m}}\,T_0$}
\opt{D}{$T=\sqrt{\dfrac{M+m}{M}}\,T_0$}
\end{qoptions}
\qn{10} Two blocks A(5kg) and B(2kg) attached to the ends of a spring of spring constant 1120 N/m are placed on a smooth horizontal plane with the spring undeformed. Simultaneously velocities of 3m/s and 10 m/s along the line of the spring in the same direction are imparted to A and B at $t=0$ then:
\qfig{p1_phy_s2_q10_1}
\begin{qoptions}
\opt{A}{When the extension of the spring is maximum the velocities of A and B are zero.}
\opt{B}{The maximum extension of the spring is 25 cm.}
\opt{C}{The first maximum compression occurs at $t=\dfrac{3\pi}{56}$ second.}
\opt{D}{Maximum extension and maximum compression occur alternately.}
\end{qoptions}
\sect{SECTION-II(i)}
\paraq{1}{2}
A projectile is thrown from a point `O' on the ground at an angle $\theta$ from horizontal with initial velocity $v_0$ at $t=0$. At the highest point of its trajectory, the projectile splits into two parts having mass in the ratio 1 : 2. After splitting, smaller part falls vertically downwards with initial velocity zero. The other part `$t$' second after splitting falls on the ground at a distance $x$ meter from the point of projection. (Consider $v_0=20$ m/s, $\theta=37^\circ$ \& acceleration due to gravity $g=10$ m/s$^2$)
\qfig{p1_phy_s3_para1_2_1}
\qn{1} The value of $t$ is\underline{\hspace{1cm}}
\qn{2} The value of $x$ is \underline{\hspace{1.2cm}}
\paraq{3}{4}
In a uniformly charged fixed nonconducting sphere, a very thin tunnel has been made along the diameter as shown in figure. A charge particle $-q$ having mass $m$ is released from rest at $t=0$ at one end of the tunnel. Total charge $Q$ of sphere is uniformly distributed in its volume. Charge particle reaches at the centre of solid sphere first time at $t=\pi\sqrt{\dfrac{\alpha.\,\pi\epsilon_0 mR^3}{Q.\,q}}$ with speed $v=\sqrt{\dfrac{Qq}{N.\,\pi\epsilon_0 mR}}$.
\qfig{p1_phy_s3_para3_4_1}
\qn{3} The value of $N$ is\underline{\hspace{1cm}}
\qn{4} The value of $\alpha$ is \underline{\hspace{1.2cm}}
\paraq{5}{6}
In figure $L$ is half part of an equiconvex lens of refractive index 1.5 whose surfaces have radius of curvature 40 cm and its right surface is silvered as shown in figure. Normal to its principal axis, a plane mirror $M$ is placed on right of lens. Distance between lens $L$ and mirror $M$ is $b$. A small object $O$ is placed on left of the lens. Images found by lens and mirror coincide with each other and tranverse length of final image formed by the lens is twice that of image formed by the mirror.
\qfig{p1_phy_s3_para5_6_1}
\qn{5} The value of $a$ (in cm) is
\qn{6} The value of $b$ (in cm) is
\sect{SECTION-II(ii)}
\qn{7} In a compound microscope the objective and eye piece have focal lengths of 0.95 cm and 5 cm respectively, and are kept at a distance of 20 cm. The last image is formed at a distance of 25 cm from eye piece. Calculate the magnitude of total magnification of the compound microscope.
\qn{8} Bottom of a glass beaker is made of a thin equi-convex lens having bottom side silvered as shown in figure. Water is filled in the beaker up to a height 4m.The image of a point object $O$, floating at middle point of beaker at the surface of water coincides with it. Find out the radius (in cm) of curvature of the lens. Given that refractive index of glass is $\dfrac32$ and that of water is $\dfrac43$.
\qfig{p1_phy_s4_q08_1}
\qn{9} A ring of radius $R$ that lies in the $yz$ plane carries a positive charge $Q$ uniformly distributed over its circumference. A particle of mass $m$ that carries a negative charge of magnitude $q$ is at the center of the ring. When the charges $Q$ and $q$ are 5 $\mu$C and $-5$ $\mu$C, respectively, and the radius of the ring is 8.0 cm, the mass $m$ oscillates about its equilibrium position with an angular frequency of 16 rad/s. Find the angular frequency (in rad/s) of oscillation of the mass if the radius of the ring is made four times to 32 cm and all other parameters remain unchanged. (Round off to nearest integer).
\raggedcolumns
\partsec{2}{CHEMISTRY}
\sect{SECTION-I(i)}
\qn{1} \qfigi{p1_chem_s1_q01_1}
\begin{qoptions}
\opt{A}{\qfigi{p1_chem_s1_q01_A}}
\opt{B}{\qfigi{p1_chem_s1_q01_B}}
\opt{C}{\qfigi{p1_chem_s1_q01_C}}
\opt{D}{\qfigi{p1_chem_s1_q01_D}}
\end{qoptions}
\qn{2} Reaction correctly matched with its major product is :
\begin{qoptions}
\opt{A}{\qfigi{p1_chem_s1_q02_A}}
\opt{B}{\qfigi{p1_chem_s1_q02_B}}
\opt{C}{\qfigi{p1_chem_s1_q02_C}}
\opt{D}{1-butene $\xrightarrow{\mathrm{HBr+H_2O_2}}$ 1-bromobutane}
\end{qoptions}
\qn{3} An ideal binary solution contains 4 mol of A and 6 mol of B. Its vapour pressure is found to be 460 mmHg. To this solution, when 10 mol of B is added, its vapour pressure becomes 480 mmHg. The vapour pressures of pure A and pure B, respectively, are
\begin{qoptions}
\opt{A}{300 mmHg, 600 mmHg}
\opt{B}{400 mmHg and 500 mmHg}
\opt{C}{300 mmHg, 500 mmHg}
\opt{D}{400 mmHg and 600 mmHg}
\end{qoptions}
\qn{4} Which of the following does not obey Sidgwick effective atomic number rule -
\begin{qoptions}
\opt{A}{$\mathrm{Co_2(CO)_8}$}
\opt{B}{$\mathrm{Fe(CO)_5}$}
\opt{C}{$\mathrm{[V(CO)_6]^-}$}
\opt{D}{$\mathrm{[Mn(CO)_6]}$}
\end{qoptions}
\sect{SECTION-I(ii)}
\qn{5} \qfigi{p1_chem_s2_q05_1}
\begin{qoptions}
\opt{A}{\qfigi{p1_chem_s2_q05_A}}
\opt{B}{\qfigi{p1_chem_s2_q05_B}}
\opt{C}{\qfigi{p1_chem_s2_q05_C}}
\opt{D}{\qfigi{p1_chem_s2_q05_D}}
\end{qoptions}
\qn{6} Out of the given options reaction(s) in which the substrate `X' can be acetaldehyde to give the final product :
\begin{qoptions}
\opt{A}{$\text{`X'}\xrightarrow[\text{(ii) NH}_4\text{Cl}]{\text{(i) CH}_3\text{MgBr}}\mathrm{CH_3{-}CH(OH){-}CH_3}$}
\opt{B}{$\text{`X'}\xrightarrow{\mathrm{PCl_5}}\mathrm{CH_3{-}CHCl_2}$}
\opt{C}{$\text{`X'}\xrightarrow[\text{(ii) HOH}]{\text{(i) Mg(Hg)}}B\xrightarrow[\Delta]{\text{Conc. H}_2\text{SO}_4}\mathrm{CH_3{-}CO{-}CH_2{-}CH_3}$}
\opt{D}{$\text{`X'}\xrightarrow[\text{(ii) Br}_2/\text{CCl}_4]{\text{(i) AgOH}}\mathrm{CH_3Br}$}
\end{qoptions}
\qn{7} Which of the following statement(s) is/are correct ?
\begin{qoptions}
\opt{A}{In fluorite structure, tetrahedral voids are occupied by anions.}
\opt{B}{In zinc-blende structure, half of tetrahedral voids are occupied by cations.}
\opt{C}{In rock-salt structure, cations occupy tetrahedral voids.}
\opt{D}{The general formula of an ionic compound crystallizing in antifluorite structure is $\mathrm{A_2B}$.}
\end{qoptions}
\qn{8} Which of the following statements are not correct ?
\begin{qoptions}
\opt{A}{For a reaction $(1/2)\mathrm{A}\to2\mathrm{B}$, the rate of disappearance of A is $-2\,d[\mathrm{A}]/dt$}
\opt{B}{For a reaction $(1/2)\mathrm{A_2}\to2\mathrm{B}$, the rate of disappearance of $\mathrm{A_2}$ is related to the rate of appearance of B through an expression $-d[\mathrm{A_2}]/dt=4\,d[\mathrm{B}]/dt$.}
\opt{C}{Reactions with small activation energies slow a lesser increase in rate constant with rise in temperature as compared to that having high activation energy.}
\opt{D}{If the half-life for a given reaction having single reactant is doubled as the initial concentration of a reactant was doubled then order of the reaction would be zero.}
\end{qoptions}
\qn{9} $\mathrm{M^{2+}+NaH_2EDTA\to[M(EDTA)]^{2-}+2H^++2Na^+}$\\
Which of the following metal ion has least tendency for the above reaction.
\begin{qoptions}
\opt{A}{$\mathrm{Cs^+}$}
\opt{B}{$\mathrm{Mg^{2+}}$}
\opt{C}{$\mathrm{Ca^{2+}}$}
\opt{D}{$\mathrm{Pb^{2+}}$}
\end{qoptions}
\qn{10} Select similarity between : $\mathrm{[Co(NH_3)_4(H_2O)Cl]Cl_2}$ and $\mathrm{[Co(NH_3)_4Cl_2]Cl.H_2O}$
\begin{qoptions}
\opt{A}{Oxidation number and coordination number of metal}
\opt{B}{$\mathrm{Cl^-}$ and $\mathrm{H_2O}$ are the ligands}
\opt{C}{Amount of precipitate with $\mathrm{AgNO_3}$ solution (by equimolar 1 litre solution of both)}
\opt{D}{Number of Co--N linkages}
\end{qoptions}
\sect{SECTION-II(i)}
\paraq{1}{2}
Consider the reactions given below and answer the following questions
\qfig{p1_chem_s3_para1_2_1}
(II) (R)--2--Chloropentane $\xrightarrow[h\nu]{\mathrm{Cl_2}}$ `Y'
\qn{1} Find out total number of monochlorinated product `X'
\qn{2} If the mixture of dichlorinated products (Y) is subjected to fractional distillation, how many fractions would be obtained?
\paraq{3}{4}
For the zero-order reaction :\\
$\mathrm{4A(g)\to3B(g)+2C(g)}$\\
At $t=0$, partial pressure of A $=4000$ mm Hg\\
Rate constant of reaction $=5$ mm Hg /sec
\qn{3} Find the total final pressure of gas mixture at $t=250$ second (in mm Hg)
\qn{4} Find the partial pressure of `B' at $t=100$ second (in mm Hg)
\paraq{5}{6}
\qfig{p1_chem_s3_q04_1}
\qn{5} The value of $\left(\dfrac{X}{Y}\right)$ is...........
\qn{6} The ratio of number of $\pi$ bond(s) to number of $\sigma$ bond(s) in `Z' is ...........
\qfullbegin
\sect{SECTION-II(ii)}
\qn{7} How many compounds out of following use two or more than two moles of Grignard reagent to convert into alcohol after complete hydrolysis.
\qfig{p1_chem_s4_q07_1}
\qfullend
\qn{8} For the reaction $\mathrm{2NO+H_2\to N_2O+H_2O}$ the value of $-dp/dt$ was found to be 1.50 Torr s$^{-1}$ for a pressure of 304 Torr of NO and 0.375 Torr s$^{-1}$ for a pressure of 152 Torr of NO, the pressure of $\mathrm{H_2}$ being constant. On the other hand, when the pressure of NO was kept constant, $-dp/dt$ was 1.60 Torr s$^{-1}$ for a hydrogen pressure of 294 Torr and 0.8 Torr s$^{-1}$ for a hydrogen pressure of 147 Torr. Determine the overall order of the reaction.
\qn{9} Select complex in which metal--carbon linkage(s) is / are present :\\
$\left[\overset{\mathrm{III}}{\mathrm{Fe}}(\mathrm{EDTA})\right]^{-}$, $\mathrm{K[PtCl_3(\pi\text{--}C_2H_4)]}$, $\mathrm{[Cr(acac)_3]^\circ}$, $\mathrm{[Fe(CO)_5]}$, $\mathrm{[Co(C_2H_5O)_6]^{3-}}$, $\mathrm{[Ni(dmg)_2]}$
\raggedcolumns
\partsec{3}{MATHEMATICS}
\sect{SECTION-I(i)}
\qn{1} If $y=\dfrac{ax+b}{cx+d}$, then value of $\dfrac{y_2^2}{y_1y_3}$, wherever exists, is (where $y_n$ denotes $n^{\mathrm{th}}$ order derivative of differentiable function $y$ and $a,b,c,d$ are constants) -
\begin{qoptions}
\opt{A}{$\dfrac32$}
\opt{B}{$\dfrac23$}
\opt{C}{$\dfrac45$}
\opt{D}{$\dfrac54$}
\end{qoptions}
\qn{2} If $f(x)=\begin{vmatrix}(x-1)^2&(x-2)^3&x+1\\(x-1)^3&x^2&(x+1)^2\\x+2&x+1&x^5\end{vmatrix}$ then coefficient of $x$ in $f(x)$ is -
\begin{qoptions}
\opt{A}{$-14$}
\opt{B}{$-16$}
\opt{C}{0}
\opt{D}{25}
\end{qoptions}
\qn{3} Let $f(x)$ be a function satisfying the condition $f(|x-3|)+2f(|x+3|)=x^2+2x\ \forall\ x\in R$, then which of the following is not correct
\begin{qoptions}
\opt{A}{minimum value of $f(x)$ is $-3$}
\opt{B}{$f(x)$ is one-one function}
\opt{C}{equation $|f(x)|=2$ has four distinct solutions}
\opt{D}{$|f(x)|$ is non-differentiable at two points}
\end{qoptions}
\qn{4} For three invertible functions, $f(x)=\dfrac{x-3}{x+3}$, $g(x)=2x$ and $h(x)=2^{\frac{x+3}{x+2}}$, value of $f^{-1}(g^{-1}(h^{-1}(4)))$ is -
\begin{qoptions}
\opt{A}{$-\dfrac12$}
\opt{B}{$\dfrac12$}
\opt{C}{1}
\opt{D}{$-1$}
\end{qoptions}
\sect{SECTION-I(ii)}
\qn{5} Given $\displaystyle\lim_{x\to\alpha}\dfrac{\sin\left(\left[x^3\right]-[x]^3\right)}{(x-\alpha)^3}$ exists finitely, then the value of $\alpha$ can be (where $\alpha\in I$ \& $[.]$ denotes greatest integer function) -
\begin{qoptions}
\opt{A}{0}
\opt{B}{1}
\opt{C}{$-1$}
\opt{D}{2}
\end{qoptions}
\qn{6} Let $f:R\to A$ be an invertible function defined by $f(x)=\log_2\left(2^{x^3+x}+4\right)$. If $g(x)$ be its inverse function, then -
\begin{qoptions}
\opt{A}{$g(3)=2$}
\opt{B}{$g'(3)=\dfrac12$}
\opt{C}{$\displaystyle\lim_{x\to3}(g(x))^{\frac{\ln4}{x-3}}=2$}
\opt{D}{set $A$ is $(2,\infty)$}
\end{qoptions}
\qn{7} If the function $f(x)=-4e^{\frac{1-x}{2}}+1+x+\dfrac{x^2}{2}+\dfrac{x^3}{3}$ and $g(x)=f^{-1}(x)$, then -
\begin{qoptions}
\opt{A}{$g'\left(-\dfrac76\right)=\dfrac15$}
\opt{B}{$g'\left(\dfrac76\right)=\dfrac15$}
\opt{C}{$g(1)=-\dfrac76$}
\opt{D}{$g\left(-\dfrac76\right)=1$}
\end{qoptions}
\qn{8} For which of the following values of `$m$' the line $y=mx$ cuts the graph of $y=\sin^{-1}(\sin x)$ in exactly 5 points -
\begin{qoptions}
\opt{A}{$\dfrac15$}
\opt{B}{$-\dfrac14$}
\opt{C}{$-\dfrac18$}
\opt{D}{$-\dfrac19$}
\end{qoptions}
\qn{9} Let $f(x)=x-1$ \& $g(x)=\dfrac1x$, then the set of points, where $[(gofog)(x)]$ is discontinuous is ($[.]$ denotes greatest integer function) -
\begin{qoptions}
\opt{A}{$\{0\}$}
\opt{B}{$\left\{\dfrac{n}{n+1}\right\}$, $n\in I-\{-1\}$}
\opt{C}{$\left\{\dfrac{n}{n+1}\right\}\cup\{1\}$, $n\in I-\{-1\}$}
\opt{D}{$\{0,1\}$}
\end{qoptions}
\qn{10} If $\sec^{-1}\left(\sqrt{1+x^2}\right)+\mathrm{cosec}^{-1}\left(\dfrac{\sqrt{1+y^2}}{y}\right)+\cot^{-1}\left(\dfrac1z\right)=\pi$ $(x,y,z>0)$ then $x+y+z$ is always -
\begin{qoptions}
\opt{A}{equal to $xyz$}
\opt{B}{greater than or equal to $3(xyz)^{1/3}$}
\opt{C}{greater than or equal to $4(xyz)^{1/3}$}
\opt{D}{greater than or equal to $5(xyz)^{1/3}$}
\end{qoptions}
\sect{SECTION-II(i)}
\paraq{1}{2}
Given $g(x)=\sin\dfrac1x$, $h(x)=x^2+2x+(\lambda+1)$ and $u(x)=\dfrac1x+\cos\dfrac1{x^2}$\\
Let $f(x)=\displaystyle\lim_{n\to\infty}\dfrac{x^{2n+1}g(x)+h(x)}{x^{2n}+3x.\,u(x)}$
\qn{1} If $\displaystyle\lim_{x\to0}f(x)=2$, then the value of $\lambda$ is -
\qn{2} $\displaystyle\lim_{x\to\infty}f(x)$ is equal to -
\paraq{3}{4}
Let $f:R\to[0,\infty)$ be a derivable function which satisfies $f(x^2+y^2)=f(2xy)+f(x^2-y^2)\ \forall\ x,y\in R$. Also $f'(1)=2$.
\qn{3} $f(3)+f(4)$ is equal to -
\qn{4} Number of solutions of $f(\cos\theta)=f(\sin\theta)$ in $[0,2\pi]$ -
\paraq{5}{6}
Consider $f(x)=x^4-3x^3+2$ and $h(x)=\dfrac{f(x)}{g(x)}$ be a rational function.\\
(i) $\displaystyle\lim_{x\to\infty}h(x)=$ infinite\\
(ii) $\displaystyle\lim_{x\to1}h(x)=\dfrac52$\\
(iii) $h(x)$ is continuous $\forall\ x\in R-\{1,2\}$\\
(iv) $g(x)$ does not have any repeated factors.
\qn{5} Number of points where $P(x)=\mathrm{sgn}(g(x))$ is non differentiable in $R$ is (where $\mathrm{sgn}(x)$ denotes signum function of $x$)
\qn{6} Number of solutions of equation $|g(|x|)|=\dfrac12$ is -
\sect{SECTION-II(ii)}
\qn{7} If $[x]^3+[y]^3+[z]^3-3[x][y][z]=4$ where $x,y,z>0$ then maximum possible value of $([x]+[y])^{[z]}$ is (where $[.]$ represents greatest integer function)
\qn{8} If a particle starts moving from $(0,0)$ and rate of change of its abscissa is 8 cm/s on the curve $y=\sqrt x$, then rate of change of area of triangle formed by co-ordinate axes \& tangent at a point reached after 2 sec of start is
\qn{9} If $y=\left(\ln\left(x+\sqrt{x^2+1}\right)\right)^2$ and $(x^2+1)\dfrac{d^3y}{dx^3}+(px+q)\dfrac{d^2y}{dx^2}+r\dfrac{dy}{dx}=sy$ , then value of $(p+q+r+s)$ is
\raggedcolumns
\newpage
\papersec{2}
\partsec{1}{PHYSICS}
\sect{SECTION-I(i)}
\qn{1} A semi-circular ring of mass $m$ and radius $R$ with linear charge density $-\lambda$, hinged at its centre with the help of a massless spoke as shown is placed in a uniform electric field as shown in the figure. If the ring is slightly rotated about $O$ in the plane of paper and released. Take $m=8$kg, $\lambda=2$C/m and $E=2$N/c.
\qfig{p2_phy_s1_q01_1}
\begin{qoptions}
\opt{A}{ring will perform angular SHM}
\opt{B}{the time period of oscillation is $2\pi$ sec.}
\opt{C}{shown position is unstable equilibrium}
\opt{D}{for large angular displacement motion will still be SHM}
\end{qoptions}
\qn{2} Two large non conducting plates having surface charge densities $+\sigma$ and $-\sigma$, respectively, are fixed `$d$' distance apart. A small test charge $q$ of mass $m$ is attached to two non-conducting identical springs of spring constant $k$ as shown in the adjacent figure. The charge $q$ is now released from rest with springs in natural length. Then $q$ will [neglect gravity] :-
\qfig{p2_phy_s1_q02_1}
\begin{qoptions}
\opt{A}{perform SHM with angular frequency $\sqrt{\dfrac{2k}{m}}$}
\opt{B}{perform SHM with amplitude $\dfrac{\sigma q}{2k\epsilon_0}$}
\opt{C}{not perform SHM, but will have a periodic motion if charges are removed on plates as well as on $m$}
\opt{D}{remain stationary}
\end{qoptions}
\qn{3} An electric dipole is placed along $x$-axis with its centre at origin :
\begin{qoptions}
\opt{A}{On a $y$-$z$ plane passing through origin, the electric field is perpendicular at each and every point on that plane at a point far from dipole.}
\opt{B}{$y$-$z$ plane passing through origin is an equipotential surface}
\opt{C}{$y$-$z$ plane passing through origin is not an equipotential surface}
\opt{D}{Flux of electric field through a closed surface enclosing dipole is zero}
\end{qoptions}
\qn{4} Block $m$ is released from rest when spring is in its natural length (assume pulley is ideal and block does not strike on ground during it's motion in vertical plane) then :
\qfig{p2_phy_s1_q04_1}
\begin{qoptions}
\opt{A}{maximum speed of block is $g\sqrt{\dfrac mk}$}
\opt{B}{maximum speed of block is $2g\sqrt{\dfrac mk}$}
\opt{C}{maximum elongation in spring is $4\,mg/k$}
\opt{D}{maximum elongation in spring is $2\,mg/k$}
\end{qoptions}
\qn{5} Two refracting media are separated by a spherical interface as shown in the figure. $PP'$ is the principal axis, $\mu_1$ and $\mu_2$ are the refractive indices of medium of incidence and medium of refraction respectively. Then :
\qfig{p2_phy_s1_q05_1}
\begin{qoptions}
\opt{A}{if $\mu_2=1.5$ and $\mu_1=1.2$, then there cannot be a real image of real object.}
\opt{B}{if $\mu_2=1.5$ and $\mu_1=1.2$, then there cannot be a real image of virtual object.}
\opt{C}{if $\mu_2=1.2$ and $\mu_1=1.5$, then there cannot be a virtual image of virtual object.}
\opt{D}{if $\mu_2=1.2$ and $\mu_1=1.5$, then there cannot be a real image of real object.}
\end{qoptions}
\qn{6} The light on reflection from a plane mirror can give a real image when :
\begin{qoptions}
\opt{A}{the convergent rays are incident on the mirror}
\opt{B}{the divergent rays are incident on the mirror}
\opt{C}{an object is placed very close to the mirror}
\opt{D}{an object is placed very far away from the mirror}
\end{qoptions}
\sect{SECTION-I(ii)}
\paraq{7}{8}
Two identical plane mirrors AB and CD of length $L$ each are placed parallel and a point source of light $O$ is placed midway between them as shown. All the shown dashed segments are either parallel or perpendicular to the mirrors. Answer the following questions based on this information.
\qfig{p2_phy_s2_para7_8_1}
\qn{7} Choose the correct statement :
\begin{qoptions}
\opt{A}{All the ray's incident on mirror AB directly from source $O$, after reflection will be incident on mirror CD}
\opt{B}{Some of the ray's incident on mirror AB directly from source $O$, after reflection will be incident on mirror CD}
\opt{C}{All the ray's incident on mirror CD directly from source $O$, after reflection will be incident on mirror AB}
\opt{D}{Some of the ray's incident on mirror CD directly from source $O$, after reflection will be incident on mirror AB}
\end{qoptions}
\qn{8} The total number of images of the object as formed by both the mirrors is :
\begin{qoptions}
\opt{A}{2}
\opt{B}{3}
\opt{C}{4}
\opt{D}{infinite}
\end{qoptions}
\paraq{9}{10}
A rectangular wire frame is kept on the principal axis of a concave mirror as shown in figure. $F$ \& $C$ are the focus \& centre of curvature of mirror respectively
\qfig{p2_phy_s2_para9_10_1}
$f$ is the focal length of mirror. A small particle is moving along the length of wire frame with constant velocity. Based on above information answer the following questions.
\qn{9} The image of the given object is best shown by
\begin{qoptions}
\opt{A}{\qfigi{p2_phy_s2_q09_A}}
\opt{B}{\qfigi{p2_phy_s2_q09_B}}
\opt{C}{\qfigi{p2_phy_s2_q09_C}}
\opt{D}{\qfigi{p2_phy_s2_q09_D}}
\end{qoptions}
\qn{10} Mark the correct statement(s) :
\begin{qoptions}
\opt{A}{The velocity of image of particle is remaining constant}
\opt{B}{The velocity of image of particle is varying}
\opt{C}{The total length of image of wire frame is $\left(\dfrac{7+\sqrt{17}}{4}\right)f$}
\opt{D}{Both (B) and (C) are correct}
\end{qoptions}
\sect{SECTION-II(i)}
\paraq{1}{2}
An infinitely long non-conducting thin rod lying on the $z$-axis carries a uniform charge density $\lambda$ $(>0)$ in the region $z>0$ and $-\lambda$ in the region $z<0$.
\qn{1} The electric field in the $x$-$y$ plane due to the charge distribution at a point which is at a distance $r$ from the origin $O$ is $\dfrac{\lambda}{n\pi\epsilon_0r}$. Find $n$.
\qn{2} The flux due to this electric field passing through the circular region : $x^2+y^2=R^2$ in $xy$ plane is $\dfrac{n\lambda R}{4\epsilon_0}$. Find $n$.
\qfig{p2_phy_s3_q02_1}
\paraq{3}{4}
A system consists of two small spheres of mass 1kg and 2kg interconnected by a weightless spring. At $t=0$, the spheres are set in motion with initial velocities as shown in figure. Position vector of centre of mass at $t=0$ is $20\,\hat j$. Initially spring is in its natural length. System moves in earth's gravitational field.
\qfig{p2_phy_s3_para3_4_1}
\qn{3} The momentum of the system at $t=2$ sec in kg m/s is $\left(x\hat i+y\hat j\right)$. Find $(x+y)(g=10$ m/s$^2)$
\qn{4} The position vector of centre of mass of the system at $t=1$ s is $\left(x\hat i+y\hat j\right)$. Find $(x+y)$.
\paraq{5}{6}
Three balls A, B and C of identical size having uniform mass distribution are placed on a smooth horizontal surface. Given that $m_A=2m$, $m_B=m$ and $m_C=m$. Ball A collides with initial velocity $V$ with B as shown in figure. (Assume all collision to be head on). The coefficient of restitution $(e)=\dfrac12$ for all collisions
\qfig{p2_phy_s3_para5_6_1}
\qn{5} The speed of A after collision of A and B for first time is $V/n$. Find $n$.
\qn{6} The speed of B after collision of B and C for first time is $V/n$. Find $n$.
\sect{SECTION-II(ii)}
\qn{7} Two fixed, equal, positive charges, each of magnitude $q=5\times10^{-5}$ C are located at points A and B separated by a distance of 6m. An equal and opposite charge moves towards them along the line COD, the perpendicular bisector of the line AB. The moving charge, when it reaches the point C at a distance of 4m from O, has a kinetic energy of 1.5 joules. Calculate the distance of the farthest point D from A (in m) which the negative charge will reach before returning towards C.
\qfig{p2_phy_s4_q07_1}
\qn{8} Two bodies of charges $-Q$ and mass $m$ each orbit a body of charge $Q$ (at rest) in a circle of radius $r$. What is the total mechanical energy of the system? If answer is $\dfrac{-3kQ^2}{nr}$, find `$n$'. (Here $k=\dfrac{1}{4\pi\epsilon_0}$) (Neglect self energy)
\qfig{p2_phy_s4_q08_1}
\qn{9} Two masses 10kg and 20 kg are connected by a massless spring. A force of 200 N acts on 20 kg mass. At the instant when the 10 kg mass has an acceleration 12 m/s$^2$ the energy stored (in J) in the spring ($k=2400$ N/m) is.
\qfig{p2_phy_s4_q09_1}
\raggedcolumns
\columnbreak
\partsec{2}{CHEMISTRY}
\sect{SECTION-I(i)}
\qn{1} Which reactions are correctly matched with their product :
\begin{qoptions}
\opt{A}{\qfigi{p2_chem_s1_q01_A}}
\opt{B}{\qfigi{p2_chem_s1_q01_2}}
\opt{C}{\qfigi{p2_chem_s1_q01_C}}
\opt{D}{$\mathrm{(CH_3COO)_2Ca}\xrightarrow{\Delta}$\qfigi{p2_chem_s1_q01_4}}
\end{qoptions}
\qn{2} Which of the following is/are correct statement regarding major product?
\qfig{p2_chem_s1_q02_1}
\begin{qoptions}
\opt{A}{Intramolecular nucleophilic substitution reaction takes place during the reaction sequence}
\opt{B}{5 membered ring is formed as major product}
\opt{C}{Grignard reagents ($\mathrm{CH_3MgBr}$) prefer to give nucleophilic addition on polar $\pi$-bond}
\opt{D}{Cyclic ether is formed as final major product}
\end{qoptions}
\qn{3} Element $X$ crystallizes in a 12 co-ordination of FCC lattice. On applying high temperature it changes to 8 co-ordination BCC lattice. If $r$ is the radius of atom $X$ in both lattices, then , INCORRECT statement(s) is/are-
\begin{qoptions}
\opt{A}{Edge length of unit cell in 12 co-ordination lattice is $2\sqrt2\,r$}
\opt{B}{Edge length of unit cell in 8 co-ordination lattice is $\dfrac{4}{\sqrt3}r$}
\opt{C}{Ratio of densities before and after applying high temperature is $2(\sqrt2)^3:(\sqrt3)^3$}
\opt{D}{Distance of $2^{\mathrm{nd}}$ nearest neighbour in 8 co-ordination BCC lattice is $2\sqrt2\,r$}
\end{qoptions}
\qn{4} Which of the following metal carbonyl(s) have more C--O bond length than that of C--O bond in $\mathrm{Fe(CO)_5}$?
\begin{qoptions}
\opt{A}{$\mathrm{[V(CO)_6]^{\ominus}}$}
\opt{B}{$\mathrm{[Ti(CO)_6]^{2\ominus}}$}
\opt{C}{$\mathrm{[Mn(CO)_6]^{\oplus}}$}
\opt{D}{CO}
\end{qoptions}
\qn{5} Identify the correct option :
\begin{qoptions}
\opt{A}{$\mathrm{Na[PtBrCl(NH_3)(NO_2)]}$\newline IUPAC-Sodium amminebromidochloridonitrito-N-platinate(I)}
\opt{B}{$\mathrm{[Co(en)_3]Cl_3}$\newline IUPAC - tris(ethane-1,2-diamine)cobalt(III) chloride.}
\opt{C}{$\mathrm{[PtCl(NH_2Me)(NH_3)_2]Cl}$\newline IUPAC -diamminechlorido(methylamine)\newline platinum(II) chloride}
\opt{D}{$\mathrm{Fe(\pi\text{-}C_5H_5)_2}$\newline IUPAC-Bis($\eta^5$-cyclopentadienyl)iron(II)}
\end{qoptions}
\qn{6} Choose correct statement/s.
\begin{qoptions}
\opt{A}{$\mathrm{[Co(NH_3)_5Br]SO_4}$ and $\mathrm{[Co(NH_3)_5SO_4]Br}$ both are ionization isomers of each other and in both complex all S-O bond length are identical}
\opt{B}{$\mathrm{[Cr(H_2O)_6]Cl_3}$ and $\mathrm{[Cr(H_2O)_4Cl_2]Cl.2H_2O}$ are hyrate isomers of each other and both isomers can be distinguished by conductivity measurements}
\opt{C}{$\mathrm{[Co(NH_3)_5NO_2]Cl_2}$ can show ionization isomerism as well as linkage isomerism}
\opt{D}{$\mathrm{[Pt(NH_3)_2Cl_2]}$ and $\mathrm{[Pt(NH_3)_4][PtCl_4]}$ are polymerization isomer of each other and both show electrical conductance in aqueous solution.}
\end{qoptions}
\qfullbegin
\sect{SECTION-I(ii)}
\paraq{7}{8}
In an E2 reaction, following one step mechanism is involved.
\qfig{p2_chem_s2_para7_8_1}
\qfullsplit
As shown in the above mechanism, a $\beta$-proton anti to leaving group is abstracted by the alkoxide base. If a less bulky base is used, $\beta$-proton is eliminated giving most substituted, most stable product as the major product. On the otherhand if a bulky base is used, $\beta$-proton is abstracted from $\beta$-carbon where there is least steric hindrance and this gives least substituted, although least stable, alkene as the major product.
\qfullend
\qn{7} What is the major product in the following elimination reaction ?
\qfig{p2_chem_s2_q07_1}
\begin{qoptions}
\opt{A}{\qfigi{p2_chem_s2_q07_A}}
\opt{B}{\qfigi{p2_chem_s2_q07_B}}
\opt{C}{\qfigi{p2_chem_s2_q07_C}}
\opt{D}{\qfigi{p2_chem_s2_q07_D}}
\end{qoptions}
\qn{8} In the following reaction, how many elimination products would be formed in principle by E2 mechanism ?
\qfig{p2_chem_s2_q08_1}
\begin{qoptions}
\opt{A}{2}
\opt{B}{3}
\opt{C}{4}
\opt{D}{5}
\end{qoptions}
\paraq{9}{10}
Two liquids A and B mix to form an ideal binary liquid solution. On adding solid non-volatile solute C to the solution, A starts polymerising into a hexamer soluble solid `$\mathrm{A_6}$' (non-electrolyte) following zero-order kinetics with a rate constant of reaction $=\left(\dfrac{10^{-1}}{6}\ \text{moles/min.}\right)$ If initially 10 moles of A and 20 moles of B were taken to form liquid solution and 5 moles of C were added, then answer the following questions based on above information and data given below :\\
\relax[Data : vapour pressure of pure A $=100$ torr, vapour pressure of pure B $=90$ torr].
\qn{9} The vapour pressure of the solution 1 hour after solid C is added will be given by :
\begin{qoptions}
\opt{A}{$\dfrac{220}{3}$ torr}
\opt{B}{98 torr}
\opt{C}{$\dfrac{2220}{29}$ torr}
\opt{D}{$\dfrac{2220}{24}$ torr}
\end{qoptions}
\qn{10} Which of the following statement in incorrect regarding the above process ?
\begin{qoptions}
\opt{A}{The vapour pressure will keep on decreasing initially on adding solute C.}
\opt{B}{The vapour pressure will become constant 100 min after addition of C.}
\opt{C}{Two hours after addition of C, the vapour pressure will become 67.5 torr.}
\opt{D}{Four hours after addition of C, the vapour pressure will become 73.75 torr.}
\end{qoptions}
\sect{SECTION-II(i)}
\paraq{1}{2}
Heating of the compounds generally loose $\mathrm{CO_2}$ and $\mathrm{H_2O}$. The process is called decarboxylation and dehydration.
\qn{1} Which of the following di-carboxylic acid on heating gives cyclic anhydride?\\
(i) Oxalic acid (ii) Succinic acid (iii) Glutaric acid (iv) Maleic acid (v) Fumaric acid
\qn{2} Which of the following compounds on heating gives cyclic ester?\\
(i) $\gamma$-hydroxy acid \quad (ii) Pimelic acid \quad (iii) Urea \quad (iv) Benzoic acid \quad (v) $\delta$-hydroxy acid
\paraq{3}{4}
Following plot is drawn for a first order reaction $\mathrm{A\to2B}$ [Given $\ln2=0.7$ ]
\qfig{p2_chem_s3_para3_4_1}
\qn{3} What is the half life of reaction in minutes ?
\qn{4} What is the initial rate of reaction (in M/min.) ?
\paraq{5}{6}
According to Werner's coordination theory every central metal atom show two type of valencies in the complexes and try to satisfy both of it's valencies.
\qn{5} For complex $\mathrm{CoCl_3.5NH_3}$\\
If P.V. $=x$\\
\hspace*{1em}S.V. $=y$\\
and total number of ions in aqueous solution per molecule of complex is $=z$\\
then calculate $\dfrac{z}{x+y}=$ ?
\qn{6} How many moles of AgCl(ppt.) is obtained when one mol aqueous solution of $\mathrm{PtCl_4.4NH_3}$ is treated with excess of $\mathrm{AgNO_3}$ solution.
\sect{SECTION-II(ii)}
\qn{7} The rate constant for decomposition of $\mathrm{COCl_2(g)}$ according to reaction :\\
$\mathrm{COCl_2(g)\rightleftharpoons CO(g)+Cl_2(g)}$ ; $\Delta H=19$ kcal mole$^{-1}$\\
is expressed as : $\ln k=15-\dfrac{11000}{T(K)}$\\
Calculate activation energy for formation of $\mathrm{COCl_2}$ from $\mathrm{CO(g)}$ \& $\mathrm{Cl_2(g)}$ in kcal mole$^{-1}$ ?
\qn{8} How many given ligands can act as $\pi$-electron donor and $\pi^*$-acceptor in complexes\\
$\mathrm{CH_2=CH_2}$, $\mathrm{NH_3}$, $\mathrm{PF_3}$, $\mathrm{C_6H_6}$, $\mathrm{C_5H_5^-}$, $\mathrm{H_2O}$
\qn{9} What is the difference between the EAN value of $\mathrm{[Fe(CO)_5]}$ \& $\mathrm{[Fe(CO)_2(NO)_2]}$
\raggedcolumns
\columnbreak
\partsec{3}{MATHEMATICS}
\sect{SECTION-I(i)}
\qn{1} If $e^{\sin\left(x^2+y^2\right)}=\tan\dfrac{y^2}{4}+\sin^{-1}x$, then $y'(0)$ can be -
\begin{qoptions}
\opt{A}{$\dfrac{1}{3\sqrt\pi}$}
\opt{B}{$-\dfrac{1}{3\sqrt\pi}$}
\opt{C}{$-\dfrac{1}{5\sqrt\pi}$}
\opt{D}{$-\dfrac{1}{3\sqrt{5\pi}}$}
\end{qoptions}
\qn{2} If $(\sin^{-1}x+\sin^{-1}w).(\sin^{-1}y+\sin^{-1}z)=\pi^2$ then $\begin{vmatrix}x^{n_1}&y^{n_2}\\z^{n_3}&w^{n_4}\end{vmatrix}$, (where $n_1,n_2,n_3,n_4\in N$) -
\begin{qoptions}
\opt{A}{has a maximum value 2}
\opt{B}{has a minimum value $-2$}
\opt{C}{is independent of $n_i$ $(i=1,2,3,4)$}
\opt{D}{can take value 0}
\end{qoptions}
\qn{3} Let $f(x+y)+f(x-y)=2f(x)\ \forall\ x,y\in R$, then which of the following is/are be correct ($f(0)=0$ \& $f$ is differentiable $\forall\ x\in R$)
\begin{qoptions}
\opt{A}{$f$ must be odd}
\opt{B}{$f$ must be even}
\opt{C}{$f'(1)=f'(2)$}
\opt{D}{$f'(2)=f'(3)=f'(4)$}
\end{qoptions}
\qn{4} Let $f:[-1,1]\to[-1,1]$, $f(x)=\sin(2\sin^{-1}x)$ \& let $g(x)=\dfrac{f(x)}{\sqrt{1-x^2}}$, then -
\begin{qoptions}
\opt{A}{$f(x)$ is one-one function}
\opt{B}{$f(x)$ is odd function}
\opt{C}{$g(x)$ is continuous in $(-1,1)$}
\opt{D}{$f(x)$ is continuous for all $x\in[-1,1]$}
\end{qoptions}
\qn{5} $\displaystyle\lim_{x\to\infty}\sin^{-1}\sin\left(\dfrac{x^2+2}{x^2+1}\right)$ is equal to (where $[.]$ and $\{.\}$ denote greatest integer and fractional part function respectively)
\begin{qoptions}
\opt{A}{$\displaystyle\lim_{x\to\infty}\left[\dfrac{x^2+2}{x^2+1}\right]$}
\opt{B}{$\displaystyle\lim_{x\to\infty}\dfrac{x^2+2}{x^2+1}$}
\opt{C}{$\displaystyle\lim_{x\to\infty}\left\{\dfrac{x^2+2}{x^2+1}\right\}$}
\opt{D}{$\displaystyle\lim_{x\to\infty}\left\{\dfrac{x^2}{x^2+1}\right\}$}
\end{qoptions}
\qn{6} Consider
\[ f(x)=\begin{cases}x^2-2,&x<0\\x,&0\leqslant x<2\\(x-2)^2,&2\leqslant x<4\\x-4,&x\geqslant4\end{cases} \]
If the equation $f(x)=k$ has four roots, then $k$ can be equal to -
\begin{qoptions}
\opt{A}{$-2$}
\opt{B}{0}
\opt{C}{1}
\opt{D}{2}
\end{qoptions}
\sect{SECTION-I(ii)}
\paraq{7}{8}
A line passes through the point $P(2,3)$ makes an angle $\theta$ with positive direction of $x$-axis and it meets the lines represented by $x^2-2xy-y^2=0$ at the points $A$ and $B$.\\
On the basis of above information, answer the following questions :
\qn{7} $\left(\dfrac{d}{d\theta}(PA+PB)\right)_{\theta=\frac\pi4}$ is equal to -
\begin{qoptions}
\opt{A}{0}
\opt{B}{$\sqrt2$}
\opt{C}{$4\sqrt2$}
\opt{D}{$8\sqrt2$}
\end{qoptions}
\qn{8} If $PA.PB=17$, then $\theta$ can be equal to -
\begin{qoptions}
\opt{A}{$\dfrac\pi3$}
\opt{B}{$\dfrac\pi6$}
\opt{C}{$\dfrac\pi2$}
\opt{D}{$\dfrac{2\pi}{3}$}
\end{qoptions}
\paraq{9}{10}
Consider the functions $f(x)=\dfrac{2^x-1}{2^x+1}$; $g(x)=\sin^{-1}x$ , $gof:R\to\left[-\dfrac\pi2,\dfrac\pi2\right]$ where $gof(x)=g(f(x))$.\\
On the basis of above information, answer the following questions :
\qn{9} $gof$ is -
\begin{qoptions}
\opt{A}{one-one onto}
\opt{B}{one-one into}
\opt{C}{many one}
\opt{D}{many one into}
\end{qoptions}
\qn{10} $\left[\displaystyle\lim_{x\to-\infty}f(g(f(x)))\right]$ is (where $[.]$ greatest integer function) -
\begin{qoptions}
\opt{A}{$-1$}
\opt{B}{$-2$}
\opt{C}{0}
\opt{D}{1}
\end{qoptions}
\sect{SECTION-II(i)}
\paraq{1}{2}
Let $f(x)=\sin^{-1}x+\sin^{-1}(1-x)-\cos^{-1}x$, then answer the following
\qn{1} Number of root(s) of the equation $f(x)=0$ is -
\qn{2} If $G(x)=f(x)+f(-x)$ then value of $G(0)$ is equal to
\paraq{3}{4}
$A(0,0)$ and $B(8,2)$ are fixed points on the curve $y^3=|x|$. A point $C$ starts moving from origin along the curve $(x<0)$ such that rate of change in the ordinate is 2cm/sec. After $t_0$ seconds triangle $ABC$ become a right triangle.
\qn{3} Rate of change in area of trinagle after $t_0$ second is $k$ then $2k$ is equal to
\qn{4} After $t_0$ seconds, tangent is drawn to the curve at point $C$ to intersect it again at $(\alpha,\beta)$ then $4\alpha+3\beta$ is -
\paraq{5}{6}
Consider $f(x)=x^4-3x^3+2$ and let $h(x)=\dfrac{f(x)}{g(x)}$ be a rational function such that\\
(i) $h(x)$ is continuous in $R-\{-1,1\}$.\\
(ii) $g(x)$ does not have any repeated linear factors.\\
(iii) As $x$ approaches infinity, $h(x)$ also approaches infinity.\\
(iv) $\displaystyle\lim_{x\to1}h(x)=\dfrac12$
\qn{5} Range of $g(x)$ is $(-\infty,k]$ then $k$ is equal to
\qn{6} $\displaystyle\lim_{x\to0}\dfrac{g(\cos x)}{5\ln\left(\sec^2x\right)}$ is equal to -
\sect{SECTION-II(ii)}
\qn{7} If $\dfrac{d^3x}{dy^3}+\left(\dfrac{dy}{dx}\right)^{-4}\left(\dfrac{d^3y}{dx^3}\right)-k\left(\dfrac{dy}{dx}\right)^{-5}\left(\dfrac{d^2y}{dx^2}\right)^2=0$ is an identity, then the value of $(k)$ is
\qn{8} If $\alpha,\beta,\gamma$ are the roots of $4x^3-16x^2+19x-6=0$, then the value of the determinant
\[ \begin{vmatrix}\alpha-\beta-\gamma&2\alpha&2\alpha\\2\beta&\beta-\gamma-\alpha&2\beta\\2\gamma&2\gamma&\gamma-\alpha-\beta\end{vmatrix} \]
is
\qn{9} Let $f(x)=g(x).\,\mathrm{sgn}(x^3+6x^2+11x+6)$, where $\mathrm{sgn}$ denotes signum function \& $g(x)=ax^3+x^2+x+1$, then minimum number of points of discontinuity of the function $f(x)$ is
\end{multicols}
\end{document}
